\documentclass[conference]{IEEEtran}
\IEEEoverridecommandlockouts
% The preceding line is only needed to identify funding in the first footnote. If that is unneeded, please comment it out.
\usepackage{cite}
\usepackage{amsmath,amssymb,amsfonts}
\usepackage{algorithmic}
\usepackage{graphicx}
\usepackage{textcomp}
\usepackage{xcolor}
\def\BibTeX{{\rm B\kern-.05em{\sc i\kern-.025em b}\kern-.08em
    T\kern-.1667em\lower.7ex\hbox{E}\kern-.125emX}}
\begin{document}

\title{Análise de Eficiência e Ocupação da Rede Elétrica de Distribuição: Impacto da Iluminação Pública LED e Integração de Cargas de Mobilidade Elétrica\\

\thanks{Identify applicable funding agency here. If none, delete this.}
}

\author{\IEEEauthorblockN{1\textsuperscript{st} Luna Gomes }
\IEEEauthorblockA{\textit{dept. de Engenharia Informática} \\
\textit{ISEP}\\
Porto, Portugal \\
1231651@isep.ipp.pt}
\and
\IEEEauthorblockN{2\textsuperscript{nd} Mariana Martins}
\IEEEauthorblockA{\textit{dept. de Engenharia Informática} \\
\textit{ISEP}\\
Porto, Portugal \\
1230679@isep.ipp.pt}
\and
\IEEEauthorblockN{3\textsuperscript{rd} Samara Miranda}
\IEEEauthorblockA{\textit{dept. de Engenharia Informática} \\
\textit{ISEP}\\
Porto, Portugal \\
1230432@isep.ipp.pt}
}

\maketitle

\begin{abstract}
Este artigo apresenta uma análise estatística da ocupação dos Postos de Transformação de Distribuição (PTD) em Portugal e da sua relação com a eficiência energética da Iluminação Pública (IP). A crescente eletrificação dos transportes coloca novos desafios à infraestrutura elétrica existente, exigindo uma gestão eficiente da capacidade disponível na rede. Neste trabalho realiza-se uma Análise Exploratória de Dados (EDA), análise inferencial e modelação preditiva, utilizando dados públicos da e-REDES relativos à iluminação pública e aos postos de transformação. Com base nestes dados, estimam-se variáveis que quantificam a potência libertada pela substituição de luminárias convencionais por tecnologia LED e avalia-se a viabilidade técnica da rede para integrar infraestruturas de carregamento de veículos elétricos (VE). Um modelo de regressão linear múltipla foi desenvolvido para explicar o nível médio de ocupação da rede em função de variáveis de infraestrutura elétrica e de eficiência tecnológica. Os resultados identificam padrões regionais de utilização da rede e evidenciam as limitações e potencialidades do processo de modernização tecnológica na integração de mobilidade elétrica.
\end{abstract}
 

\begin{IEEEkeywords}
Análise de Dados, Iluminação Pública, Regressão Linear, Eficiência Energética, Mobilidade Elétrica, Redes de Distribuição, VIF, ANOVA.
\end{IEEEkeywords}

\section{Introdução}
 
A transição energética e a descarbonização das cidades são atualmente objetivos centrais das políticas energéticas europeias. A eletrificação de setores como os transportes e a crescente integração de tecnologias sustentáveis colocam novos desafios à gestão da infraestrutura elétrica existente.
 
Os Postos de Transformação de Distribuição (PTD) constituem elementos críticos da rede, responsáveis por converter e distribuir energia elétrica em baixa tensão. O aumento da procura energética, impulsionado pela expansão da mobilidade elétrica, pode conduzir a situações de elevada ocupação destes equipamentos.
 
Paralelamente, a Iluminação Pública (IP) representa uma parcela significativa do consumo energético municipal. A transição progressiva das tecnologias convencionais de iluminação --- como lâmpadas de mercúrio ou sódio --- para tecnologia LED tem demonstrado um potencial disruptivo na libertação de carga na rede elétrica de baixa tensão~\cite{b1}.
 
Surge assim uma oportunidade de sinergia: a \emph{potência libertada} pela eficiência na IP pode ser o catalisador para a infraestrutura de VE, evitando investimentos massivos em novos PTDs~\cite{b4}. Este trabalho, desenvolvido no âmbito da unidade curricular de Análise de Dados em Informática (ANADI), tem como objetivo analisar a relação entre a eficiência da iluminação pública e o nível de ocupação da rede elétrica de distribuição, construindo um modelo preditivo e avaliando a viabilidade técnica de integração de carregadores de veículos elétricos de 22~kW.
 
\section{Objetivos}
O presente trabalho persegue os seguintes objetivos específicos:
\begin{itemize}
    \item \textbf{Exploração e Tratamento:} Efetuar a limpeza e integração de dados de IP e PTD, tratando valores omissos e normalizando escalas de utilização.
    \item \textbf{Caracterização Tecnológica:} Analisar o mix tecnológico da iluminação pública nacional e identificar o potencial de ganho energético por concelho.
    \item \textbf{Diagnóstico de Saturação:} Caracterizar a ocupação dos transformadores por distrito e identificar pontos críticos de sobrecarga.
    \item \textbf{Validação Estatística:} Testar se a rede possui folga significativa para a integração de cargas VE de 22~kW.
    \item \textbf{Modelação Preditiva:} Desenvolver um modelo de regressão linear múltipla para estimar a ocupação da rede e validar a sua precisão em concelhos específicos (distrito de Aveiro).
\end{itemize}

\section{Metodologia e Dados}

A análise segue as etapas de limpeza e transformação (\textit{ETL}), exploração estatística, inferência e modelação~\cite{b2,b3}, recorrendo às bibliotecas Python \texttt{Pandas}, \texttt{NumPy}, \texttt{Matplotlib}, \texttt{Seaborn}, \texttt{SciPy} e \texttt{Statsmodels}. O dashboard interativo foi desenvolvido com \texttt{Streamlit} e \texttt{Plotly}. Todos os testes de hipótese foram realizados com nível de significância $\alpha = 0.05$.
 
\subsection{Fontes de Dados}

Foram integrados dois conjuntos de dados públicos disponibilizados pela e-REDES~\cite{b4}:

\begin{enumerate}
    \item \textbf{IP\_data}: Inventário de iluminação pública por freguesia (9\,300 registos, 13 colunas), com informação sobre tipo de lâmpada, número de luminárias e potência instalada.
    \item \textbf{PTD\_data}: Dados técnicos dos postos de transformação (72\,027 registos, 11 colunas), incluindo potência instalada, nível de utilização, número de clientes e coordenadas geográficas.
\end{enumerate}

\subsection{Resumo das Técnicas Estatísticas}
 
As técnicas estatísticas aplicadas ao longo do trabalho foram:
\begin{itemize}
    \item \textbf{Estatística Descritiva}: média, quartis, desvio padrão, assimetria e curtose.
    \item \textbf{Teste de Normalidade (Shapiro-Wilk)}: aplicado como pré-condição antes de cada teste paramétrico.
    \item \textbf{Testes de Hipóteses}: t-test para uma amostra, t-test de Welch para duas amostras independentes e ANOVA de uma via com análise post-hoc de Tukey.
    \item \textbf{Correlação de Pearson}: quantificação da intensidade e significância da relação linear entre variáveis de infraestrutura.
    \item \textbf{Regressão Linear Múltipla (OLS)}: modelação preditiva do nível de ocupação, com verificação de normalidade, independência e homocedasticidade dos resíduos, e cálculo do VIF para diagnóstico de multicolinearidade.
\end{itemize}
 

\section{Tratamento e Engenharia de Variáveis (ETL)}
\subsection{Processamento da Iluminação Pública}
 
No dataset IP\_data, criou-se a variável binária \texttt{Is\_Ineficiente} (1 para Sódio ou Mercúrio; 0 para outros tipos) e converteu-se a potência para kW. Os dados foram agregados por \texttt{CodDistritoConcelho} para obter:
\begin{align}
P\_IP\_Total_c &= \textstyle\sum_{i \in c} P_i^{(kW)} \\
P\_IP\_Inef_c &= \textstyle\sum_{\substack{i \in c,\\ \text{Is\_Inef}=1}} P_i^{(kW)}
\end{align}

\subsection{Processamento dos PTDs}
 
O campo ``Nível de Utilização [\%]'' foi normalizado para decimal extraindo o limite superior de cada intervalo (e.g., ``60\%-79\%'' $\to$ 0.79). Os dados foram agregados por concelho para obter a capacidade total (\texttt{Cap\_PTD}), utilização média (\texttt{Util\_Media}) e número de postos (\texttt{N\_PTDs}). Identificaram-se 3\,064 valores nulos e 14\,338 valores indeterminados, perfazendo 17\,402 registos afetados (24,16\% do total).

\subsection{Variáveis de Viabilidade do Dataset Consolidado}
 
A fusão dos dois datasets via \texttt{CodDistritoConcelho} resultou em 278 concelhos. Foram criadas as seguintes variáveis de análise:
\begin{align}
\Delta P_{LED}  &= P\_IP\_Inef \times 0.65 \label{eq:dpled}\\
P_{Folga}       &= (Cap_{PTD} \times 0.92) \times (1 - Util_{Media}) \label{eq:pfolga}\\
P_{VE}          &= N_{PTDs} \times 22 \times 0.60 \label{eq:pve}\\
D               &= P_{Folga} + \Delta P_{LED} - P_{VE} \label{eq:d}\\
Rate\_Inef      &= P\_IP\_Inef \;/\; P\_IP\_Total \label{eq:rate}
\end{align}
 
O saldo $D > 0$ determina a viabilidade técnica do concelho para suportar novas cargas de VE sem comprometer a operação segura da rede.
 
\section{Análise Exploratória de Dados}
 
\subsection{Mix Tecnológico e Concentração de Potência}
 
A análise do mix tecnológico da IP revelou que uma proporção relevante da potência instalada ainda provém de tecnologias convencionais (Sódio e Mercúrio). O gráfico de barras da potência ineficiente por município evidencia que esta se concentra num grupo restrito de concelhos, essencialmente capitais de distrito e municípios de maior dimensão, definindo prioridades geográficas claras para as políticas de modernização.

\subsection{Distribuição da Ocupação por Distrito}
 
Os diagramas de extremos e quartis dos PTDs nos distritos de Aveiro, Lisboa, Porto e Setúbal evidenciaram padrões distintos. Lisboa apresenta a maior variabilidade de carga, reflexo da heterogeneidade urbana e industrial, com vários municípios próximos dos limites de capacidade. Os distritos do Norte apresentam distribuições mais compactas e homogéneas.
 
\subsection{Valores Omissos e Outliers}
 
A prevalência de dados omissos foi quantificada em 17\,402 registos (24,16\% do total de PTDs). O boxplot do nível de utilização individual mostra que 50\% dos PTDs operam entre 38\% e 58\% de capacidade. Contudo, identificaram-se outliers críticos com valores $\geq 100\%$, revelando transformadores em sobrecarga. Estes constituem estrangulamentos da rede onde a instalação de carregadores VE é tecnicamente inviável sem intervenção prévia de reforço.
 
\subsection{Estatística Descritiva por Concelho}
 
A Tabela~\ref{tab:stats_concelhos} apresenta a estatística descritiva do nível de utilização individual dos PTDs para os concelhos de Braga, Coimbra, Faro e Évora (4 casas decimais).
 
\begin{table}[htbp]
\caption{Estatística Descritiva do Nível de Utilização por Concelho}
\label{tab:stats_concelhos}
\centering
\setlength{\tabcolsep}{3.5pt}
\begin{tabular}{lccccccc}
\hline
\textbf{Concelho} & \textbf{Média} & \textbf{DP} & \textbf{Q1} & \textbf{Mediana} & \textbf{Q3} & \textbf{Assim.} & \textbf{Kurt.} \\
\hline
Braga   & 0.5423 & 0.2401 & 0.39 & 0.59 & 0.79 & 0.4280  & $-$0.7116 \\
Coimbra & 0.5406 & 0.2322 & 0.39 & 0.59 & 0.79 & 0.2258  & $-$0.7060 \\
Faro    & 0.5549 & 0.2126 & 0.39 & 0.59 & 0.59 & 0.3009  & $-$0.4757 \\
Évora   & 0.4546 & 0.2425 & 0.19 & 0.39 & 0.59 & 0.6680  & $-$0.4254 \\
\hline
\end{tabular}
\end{table}
 
Faro é o concelho com a rede mais carregada em média (55,49\%), enquanto Évora apresenta a maior folga (45,46\%). A assimetria positiva em todos os concelhos indica que existem mais PTDs em cargas baixas/médias do que em sobrecarga extrema. Os valores de curtose negativos (distribuições platicúrticas) revelam que as cargas estão amplamente distribuídas pelos vários níveis de utilização.

\section{Inferência Estatística}
 
\subsection{Teste de Hipótese: Ocupação Média Inferior a 60\%}
 
Selecionou-se aleatoriamente uma amostra de 50 concelhos (\texttt{random\_state=42}). O teste de Shapiro-Wilk confirmou a normalidade ($p = 0.0717 > 0.05$), justificando o uso do t-test paramétrico. O teste unilateral formulado foi:
\[
H_0{:}\; \mu \geq 0.60 \quad \text{vs} \quad H_1{:}\; \mu < 0.60
\]
 
O t-test para uma amostra resultou em $p \approx 0 \ll 0.05$. Rejeita-se $H_0$: há evidências estatísticas significativas de que o nível médio de ocupação é inferior a 60\%, confirmando a existência de folga global de capacidade nos PTDs.
 
\subsection{Teste de Duas Amostras: Modernizados vs. Ineficientes}
 
Com base na mediana do rácio de LED, definiram-se dois grupos: ``Modernizados'' e ``Ineficientes''. Selecionaram-se 30 concelhos de cada grupo. O t-test de Welch para amostras independentes revelou que o estado médio de ocupação da rede não difere significativamente entre os dois grupos ($p > 0.05$), sugerindo que o nível de modernização da IP não é determinante para a carga global dos PTDs.
 
\subsection{ANOVA: Diferenças Regionais na Ocupação}
 
Constituíram-se três amostras de 25 concelhos: (1)~Norte/Centro Litoral (Porto, Braga, Coimbra), (2)~Lisboa e Litoral Sul (Lisboa, Setúbal, Aveiro) e (3)~Interior/Alentejo (Évora, Beja, Portalegre). A ANOVA de uma via revelou diferenças estatisticamente significativas entre as regiões ($p < 0.05$). A análise post-hoc de Tukey identificou que o Interior/Alentejo apresenta ocupações médias significativamente inferiores às regiões litorais. Entre as regiões litorais não se detetaram diferenças significativas.
 
\subsection{Correlação de Pearson: Cap\_PTD vs. P\_IP\_Total}
 
Para a totalidade dos 278 concelhos, o coeficiente de Pearson entre \texttt{Cap\_PTD} e \texttt{P\_IP\_Total} é:
\[
r = 0.9111, \quad p \approx 3.13 \times 10^{-108}
\]
 
O resultado evidencia uma correlação linear positiva muito forte e estatisticamente significativa ($p \ll 0.001$). O dimensionamento da infraestrutura de transformação acompanha de forma altamente proporcional a carga da IP, traduzindo a coevolução histórica destas infraestruturas municipais.
 
\section{Correlação e Regressão --- Modelo Preditivo}
 
\subsection{Tabela de Correlação nos 4 Distritos}
 
Para os distritos de Aveiro, Braga, Lisboa e Porto, a correlação de Pearson entre \texttt{Cap\_PTD} e \texttt{P\_IP\_Total} é $r = 0.9129$, confirmando a relação muito forte também nos grandes centros urbanos.
 
\subsection{Modelo de Regressão Linear Múltipla}
 
Para Portugal Continental (278 observações), ajustou-se o modelo OLS com $Y = \texttt{Util\_Media}$, $X_1 = \texttt{P\_IP\_Total}$, $X_2 = \texttt{Cap\_PTD}$ e $X_3 = \texttt{Rate\_Ineficiencia}$:
 
\begin{equation}
\hat{Y} = 0.4989 + 4.594 \!\times\! 10^{-5} X_1 - 2.7 \!\times\! 10^{-7} X_2 + 0.0127 X_3
\label{eq:modelo}
\end{equation}
 
Os principais indicadores de qualidade estão na Tabela~\ref{tab:ols}.
 
\begin{table}[htbp]
\caption{Sumário do Modelo OLS}
\label{tab:ols}
\centering
\begin{tabular}{lc}
\hline
\textbf{Estatística}  & \textbf{Valor}  \\
\hline
$R^2$                 & 0.074           \\
$R^2$ Ajustado        & 0.064           \\
F-statistic           & 7.300           \\
Prob (F)              & 0.000100        \\
N.º Observações       & 278             \\
AIC                   & $-$726.2        \\
\hline
\end{tabular}
\end{table}
 
\subsection{Verificação dos Pressupostos dos Resíduos}
 
\begin{itemize}
    \item \textbf{Normalidade (Shapiro-Wilk)}: $p = 0.0301 < 0.05$ --- pressuposto violado.
    \item \textbf{Independência (Durbin-Watson)}: estatística $= 1.433$, marginalmente abaixo do limiar de 1.5, indicando autocorrelação positiva ligeira --- pressuposto parcialmente violado.
    \item \textbf{Homocedasticidade (Breusch-Pagan)}: $p = 0.1710 > 0.05$ --- variância constante confirmada.
\end{itemize}
 
As violações da normalidade e independência são coerentes com o baixo $R^2$, refletindo a omissão de variáveis explicativas fundamentais (consumo doméstico e industrial).
 
\subsection{Análise de Multicolinearidade (VIF)}
 
\begin{table}[htbp]
\caption{Variance Inflation Factor (VIF)}
\label{tab:vif}
\centering
\begin{tabular}{lc}
\hline
\textbf{Variável}               & \textbf{VIF} \\
\hline
$X_1$ --- P\_IP\_Total          & 8.45 \\
$X_2$ --- Cap\_PTD              & 7.26 \\
$X_3$ --- Rate\_Ineficiencia    & 1.52 \\
\hline
\end{tabular}
\end{table}
 
$X_1$ e $X_2$ apresentam VIF $>$ 5, confirmando multicolinearidade severa, consistente com $r = 0.91$. A inflação da variância compromete a interpretabilidade individual dos coeficientes. $X_3$ (VIF = 1.52) está isenta de problemas de colinearidade.
 
\subsection{Comentário ao Modelo}
 
Com $R^2_{adj} = 0.064$, o modelo explica apenas 6.4\% da variabilidade da ocupação da rede. Os preditores $X_1$ e $X_2$ são estatisticamente significativos ($p < 0.001$), enquanto $X_3$ não o é ($p = 0.577$). A forte multicolinearidade entre $X_1$ e $X_2$ dificulta a separação dos efeitos individuais. O modelo funciona como indicador geral de nível de carga esperado, mas é insuficiente para capturar a complexidade real dos consumos locais.
 
\subsection{Previsões para o Distrito de Aveiro}
 
\begin{table}[htbp]
\caption{Previsões vs. Valores Reais --- Distrito de Aveiro}
\label{tab:previsoes}
\centering
\begin{tabular}{lccc}
\hline
\textbf{Concelho}  & \textbf{Y Real} & \textbf{Y Prev.} & \textbf{Erro} \\
\hline
Albergaria-a-Velha & 0.4699 & 0.5057 & $-$0.0358 \\
Anadia             & 0.5430 & 0.5155 & $+$0.0275 \\
Arouca             & 0.5274 & 0.5170 & $+$0.0104 \\
Aveiro             & 0.4755 & 0.4958 & $-$0.0203 \\
Castelo de Paiva   & 0.5432 & 0.5144 & $+$0.0288 \\
Espinho            & 0.4434 & 0.5014 & $-$0.0580 \\
Estarreja          & 0.5076 & 0.5074 & $+$0.0002 \\
\hline
\end{tabular}
\end{table}
 
As previsões concentram-se numa faixa estreita (49\%--53\%), sendo incapazes de captar a variação real. Os erros mais elevados verificam-se em Espinho ($-$5.8\%) e Albergaria-a-Velha ($-$3.6\%), corroborando o baixo poder preditivo do modelo.
 
\subsection{Impacto da Redução de 20\% na Taxa de Ineficiência}
 
Aplicando $\hat{\beta}_3 = 0.0127$ do modelo~\eqref{eq:modelo}:
\[
\Delta Y = \hat{\beta}_3 \times (-0.20) = 0.0127 \times (-0.20) = -0.00254
\]
 
A redução prevista seria de $\approx$0.25\%. Como $X_3$ não é estatisticamente significativo ($p = 0.577$), $\hat{\beta}_3$ não é estatisticamente diferente de zero. A implementação de tecnologia LED não tem impacto mensurável na desocupação global dos PTDs a nível concelhio.
 
\subsection{Identificação de Concelhos Viáveis para VE}
 
\begin{table}[htbp]
\caption{Concelhos Mais Viáveis para Carregadores VE de 22~kW}
\label{tab:viabilidade}
\centering
\begin{tabular}{lcc}
\hline
\textbf{Concelho}  & \textbf{Y Real} & \textbf{Lim. Sup. 95\%} \\
\hline
Porto              & 0.4219 & 0.5399 \\
Vila Nova de Gaia  & 0.4179 & 0.5412 \\
Lisboa             & 0.4005 & 0.5972 \\
Évora              & 0.4546 & 0.6173 \\
Beja               & 0.4467 & 0.6213 \\
\hline
\end{tabular}
\end{table}
 
A decisão de alocação de nova carga deve basear-se no limite superior do intervalo de confiança (cenário mais conservador). Porto e Vila Nova de Gaia apresentam os limites mais baixos (53.99\% e 54.12\%), com folga superior a 40\% no pior cenário previsto. Lisboa segue com 59.72\%. Dado o baixo $R^2_{adj}$, estas previsões devem ser validadas com medições locais ao nível de cada PTD antes de qualquer decisão de investimento.
 
\section{Dashboard Interativo}
 
Foi desenvolvido um dashboard interativo em \texttt{Streamlit} que integra os dois datasets e disponibiliza:
\begin{itemize}
    \item \textbf{Perfis horários} de consumo da IP (antes vs. depois da migração LED) em Plotly;
    \item \textbf{KPIs} em tempo real: potência libertada ($\Delta P_{LED}$), carga VE total ($P_{VE}$) e viabilidade global;
    \item \textbf{Simulação de cenários} via sliders: eficiência LED (0--100\%), penetração VE (0--100\%) e potência do carregador (3.7, 7.4, 11, 22~kW);
    \item \textbf{Mapa geográfico} com localização dos PTDs e saldo de viabilidade por concelho;
    \item \textbf{Tabela interativa} com capacidade instalada, folga e saldo final por concelho.
\end{itemize}
 
\section{Conclusão}
 
Este trabalho analisou a relação entre a eficiência tecnológica da IP e a capacidade de carga dos PTDs em Portugal Continental, com foco na integração de VE. A análise exploratória revelou que 24.16\% dos registos PTD contêm dados omissos ou indeterminados. O t-test confirmou folga global de capacidade nos PTDs ($p \approx 0$). A correlação de Pearson ($r = 0.9111$) evidenciou uma relação linear muito forte entre a capacidade de transformação e a carga de IP. A ANOVA confirmou diferenças regionais significativas, com o Interior/Alentejo a apresentar ocupações mais baixas. O teste de duas amostras não revelou diferença significativa entre concelhos modernizados e ineficientes, sugerindo que a modernização da IP não é o fator dominante na carga dos PTDs.
 
O modelo de regressão apresenta poder explicativo limitado ($R^2_{adj} = 6.4\%$), com multicolinearidade severa entre $X_1$ e $X_2$ e $X_3$ estatisticamente não significativo. Porto, Vila Nova de Gaia e Lisboa destacam-se como os municípios mais aptos para receber carregadores de 22~kW. Como trabalho futuro, propõe-se a inclusão de variáveis de consumo doméstico e industrial e a aplicação de modelos não lineares para melhorar o poder preditivo.
 
\section*{Agradecimentos}

Este trabalho foi realizado com o apoio de ferramentas de Inteligência Artificial Generativa (Gemini 2.5 Flash) para auxílio na estruturação de código Python e refinamento do texto em LaTeX, seguindo as normas éticas de utilização de IA da instituição.


\begin{thebibliography}{00}
\bibitem{b1} C. Heumann e M. Schomaker Shalabh, \textit{Introduction to Statistics and Data Analysis},
Springer International Publishing, 2016.
\bibitem{b2} D.\,C. Montgomery, \textit{Design and Analysis of Experiments}, 8.ª~ed.,
John Wiley \& Sons, Nova Iorque, 2013.
\bibitem{b3} W. McKinney, \textit{Python for Data Analysis: Data Wrangling with Pandas, NumPy, and Jupyter},
3.ª~ed., O'Reilly Media, 2022.
\bibitem{b4} e-REDES, ``Portal de Dados Abertos,'' [Online].
Acedido em: março de 2026.
\bibitem{b5} J.\,H. McDonald, \textit{Handbook of Biological Statistics}, 3.ª~ed.,
Sparky House Publishing, Baltimore, 2014.
\bibitem{b6} G. James, D. Witten, T. Hastie e R. Tibshirani,
\textit{An Introduction to Statistical Learning}, 2.ª~ed., Springer, 2021.
\end{thebibliography}

\end{document}
